% JAMT starter. Fill yellow fields; keep role commands and class files.
% All table/formula values below are layout examples, not research results.
\documentclass{jamt}
\begin{document}
\JAMTParagraph{article_type}{\JAMTMissing{Article type / Тип статьи}}
\JAMTParagraph{rubric}{\JAMTMissing{Journal section / Рубрика}}
\JAMTParagraph[after=4]{editorial_metadata}{\JAMTMissing{Year, volume, issue, pages / Год, том, выпуск, страницы}}
\JAMTParagraph[after=4]{editorial_metadata}{DOI: \JAMTMissing{Insert the assigned DOI}}
\JAMTParagraph[after=4]{editorial_metadata}{UDC: \JAMTMissing{Insert the UDC index}}

\JAMTParagraph{title}{\JAMTMissing{Article title in English}}
\JAMTParagraph{author}{\JAMTMissing{Authors in English}}
\JAMTParagraph{affiliation}{\JAMTMissing{Institutions, city, country in English}}
\JAMTParagraph{email}{\JAMTMissing{Corresponding author email}}
\JAMTParagraph{abstract}{\textbf{Abstract.} \JAMTMissing{Insert the original abstract in English}}
\JAMTParagraph{keywords}{\textbf{Keywords:} \JAMTMissing{Insert keywords in English}}
\JAMTParagraph{citation}{\textbf{For citation:} \JAMTMissing{Authors, title, year, volume, issue, pages, DOI}}
\JAMTFrontDivider

\JAMTParagraph[language=russian]{title}{\JAMTMissing{Название статьи на русском языке}}
\JAMTParagraph[language=russian]{author}{\JAMTMissing{Авторы на русском языке}}
\JAMTParagraph[language=russian]{affiliation}{\JAMTMissing{Организации, город, страна на русском языке}}
\JAMTParagraph{email}{\JAMTMissing{Email автора для переписки}}
\JAMTParagraph[language=russian]{abstract}{\textbf{Аннотация.} \JAMTMissing{Добавьте авторскую аннотацию на русском языке}}
\JAMTParagraph[language=russian]{keywords}{\textbf{Ключевые слова:} \JAMTMissing{Добавьте ключевые слова на русском языке}}
\JAMTParagraph[language=russian]{citation}{\textbf{Для цитирования:} \JAMTMissing{Авторы, название, год, том, выпуск, страницы, DOI}}

\begin{JAMTBody}
\JAMTHeading{1. Introduction}
\JAMTParagraph[language=russian]{body}{\JAMTMissing{Добавьте исходный текст введения}}
\JAMTParagraph[language=russian]{body}{Абзацы этого макета показывают расположение элементов статьи. Замените их содержимым рукописи. Шрифт, отступы и интервалы задаются общим профилем. Сведения, которых нет в исходнике, требуется предоставить автору или редакции.}
\JAMTHeading{2. Materials and methods}
\JAMTParagraph[language=russian]{body}{\JAMTMissing{Добавьте материалы и методы исследования}}
\JAMTParagraph[language=russian]{body}{Пример структурной формулы для проверки вёрстки, который следует заменить формулой из рукописи:}
\begin{equation}
  \bar{x}=\frac{1}{n}\sum_{i=1}^{n}x_i .
\end{equation}
\JAMTHeading{3. Results and discussion}
\JAMTParagraph[language=russian]{body}{\JAMTMissing{Добавьте результаты и обсуждение}}
\begin{JAMTObject}
\JAMTParagraph{table_caption}{Table 1. Layout example}
\JAMTParagraph[language=russian]{table_caption}{Таблица 1. Пример оформления}
{\fontsize{\labtablesize}{\labtableleading}\selectfont
\begin{tabular}{@{}p{.54\linewidth}p{.34\linewidth}@{}}
\toprule
Параметр & Значение \\
\midrule
\JAMTMissing{Параметр A} & \JAMTMissing{Данные} \\
\JAMTMissing{Параметр B} & \JAMTMissing{Данные} \\
\bottomrule
\end{tabular}}
\end{JAMTObject}
\JAMTHeading{4. Conclusions}
\JAMTParagraph[language=russian]{body}{\JAMTMissing{Добавьте выводы без изменения авторского смысла}}
\JAMTHeading[references_heading]{References / Список литературы}
\JAMTParagraph[language=russian]{reference_item}{1. \JAMTMissing{Добавьте полное описание источника из рукописи}}
\end{JAMTBody}

\JAMTHeading[author_information]{Information about the authors / Информация об авторах}
\JAMTParagraph[language=russian]{author_bio}{\JAMTMissing{Сведения об авторах на русском и английском языках}}
\JAMTParagraph{received_metadata}{Received: \JAMTMissing{Date supplied by the editorial office}}
\JAMTParagraph{received_metadata}{Accepted: \JAMTMissing{Date supplied by the editorial office}}
\JAMTParagraph{received_metadata}{Published: \JAMTMissing{Date supplied by the editorial office}}
\end{document}
